\documentclass[10pt]{article}
\usepackage[margin=0.9in]{geometry}
\usepackage{amsmath,amssymb,amsthm,hyperref}
\newtheorem{theorem}{Theorem}
\newtheorem{lemma}{Lemma}
\newcommand{\R}{\mathbb R}
\title{Diagonal maximality of the inverse geometric Vandermonde matrix}
\author{Principia Math in Progress}
\date{12 September 2026}
\begin{document}
\maketitle
\paragraph{Task 1. Solve the inverse Vandermonde conjecture.}
For \(b>1\) and \(n\ge1\), put
\[
 V=(b^{ij})_{0\le i,j<n},\qquad C=V^{-1}=(c_{ij}),\qquad
 M_b(n)=\max_{i,j}|c_{ij}|,\qquad
 n_0=\left\lceil\log_b(1+b^{-1})\right\rceil.
\]
Sanna, Shallit, and Zhang conjectured that, for all sufficiently large \(n\),
\(M_b(n)=|c_{ii}|\) for some \(0\le i\le n_0\)
\cite[Conjecture 11]{SSZ}. The accepted manuscript retains this as
Conjecture 12. We prove the assertion for every \(n\).

\begin{theorem}
For every \(b>1\) and \(n\ge1\),
\[
 M_b(n)=\max_{0\le i\le\min(n-1,n_0)}c_{ii}.
\]
All diagonal entries are positive, and no off-diagonal entry attains \(M_b(n)\).
\end{theorem}

\begin{proof}
Write \(s=\log b>0\). For \(x=(x_0,\ldots,x_{n-1})\in\R^n\),
the absolutely convergent exponential series gives
\[
 x^TVx=\sum_{k=0}^{\infty}\frac{s^k}{k!}
             \left(\sum_{i=0}^{n-1}x_i i^k\right)^2 ,
 \qquad (0^0=1).
\]
If \(x\ne0\), at least one inner sum with \(0\le k<n\) is nonzero:
the Vandermonde matrix on the distinct nodes \(0,1,\ldots,n-1\)
is invertible. Thus \(V\), and hence \(C\), is positive definite.
In particular \(c_{ii}>0\), and each two-by-two principal minor yields
\[
 |c_{ij}|^2<c_{ii}c_{jj}\quad(i\ne j).
 \tag{1}
\]
Consequently \(M_b(n)=\max_i c_{ii}\), and off-diagonal entries are
strictly smaller than that maximum.

For completeness we verify the needed index bound, also obtainable
from \cite[Theorem 1]{SSZ}. Let \(e_m\) denote the elementary symmetric
polynomial of degree \(m\), with \(e_0=1\), and set
\[
 E_m(j)=e_m(\{b^h:0\le h<n,\ h\ne j\}),\qquad
 \pi_j=\prod_{\substack{0\le h<n\\h\ne j}}|b^j-b^h|.
\]
The coefficients of the Lagrange polynomial
\(\prod_{h\ne j}(X-b^h)/(b^j-b^h)\) form column \(j\) of \(C\).
Thus
\[
 |c_{ij}|=\frac{E_{n-1-i}(j)}{\pi_j}.                 \tag{2}
\]
For \(0\le j<n-1\), let \(T=\{b^h:h\ne j,j+1\}\), and use
\(e_m(T)=0\) outside its natural degree range. Then
\[
 E_m(j)-E_m(j+1)=(b^{j+1}-b^j)e_{m-1}(T)\ge0.       \tag{3}
\]
Direct cancellation in the two products for \(\pi_j\) gives
\[
 \frac{\pi_{j+1}}{\pi_j}
   =\frac{b^{n+j-1}-b^{n-2}}{b^{n-1}-b^j}>1
       \quad\hbox{if } j\ge n_0.                  \tag{4}
\]
Indeed the last inequality is equivalent to
\[
 b^j>\frac{b^{n-1}+b^{n-2}}{b^{n-1}+1},
\]
whose right side is strictly less than \(1+b^{-1}\le b^j\).
All numerators in (2) are positive. Equations (2)--(4) therefore show
\[
 |c_{i,j+1}|<|c_{ij}| \quad(0\le i<n,\ n_0\le j<n-1).
\]
Using symmetry, for \(n_0\le j<n-1\) we obtain
\[
 c_{j+1,j+1}<|c_{j+1,j}|=|c_{j,j+1}|<c_{jj}.
\]
Hence the diagonal strictly decreases after index \(n_0\).
If \(n-1\le n_0\), the index restriction is automatic; otherwise
every larger diagonal entry is bounded by \(c_{n_0n_0}\).
Together with (1), this proves the theorem.
\end{proof}

\paragraph{Task 2. Locate the maximizing diagonal index.}
Determine, as functions of \(b\) and \(n\), which indices attain
\(\max_i c_{ii}\), and characterize the values of \(b\) at which ties occur.

\paragraph{Task 3. Quantify convergence to the limiting maximum.}
Find effective error bounds for \(M_b(n)\) relative to its known limit,
with explicit dependence on \(b\), especially as \(b\downarrow1\).

\paragraph{Task 4. Study other exponential kernels.}
For distinct real nodes \(t_0,\ldots,t_{n-1}\), analyze which diagonal
entries maximize the inverse of \((e^{s t_i t_j})\), and identify node
conditions that yield a dimension-independent range of maximizing indices.

\paragraph{Task 5. Measure separation from off-diagonal entries.}
Obtain explicit estimates for the gap between \(M_b(n)\) and
\(\max_{i\ne j}|c_{ij}|\), and determine its behavior for fixed \(b\)
as \(n\) increases.

\begin{thebibliography}{1}
\bibitem{SSZ}
C.~Sanna, J.~Shallit, and S.~Zhang,
\emph{The Largest Entry in the Inverse of a Vandermonde Matrix},
\href{https://arxiv.org/abs/2008.01012}{arXiv:2008.01012v1} (2020).
Accepted manuscript, 11 December 2020:
\href{https://iris.polito.it/handle/11583/2898892}{institutional record}.
Published in \emph{Linear and Multilinear Algebra} \textbf{70} (2022),
5634--5641, DOI \href{https://doi.org/10.1080/03081087.2021.1922337}
{10.1080/03081087.2021.1922337}.
\end{thebibliography}
\end{document}
